\documentclass[10pt,a4paper]{amsart}
\usepackage[margin=19mm]{geometry}
\usepackage{amsmath,amssymb,mathtools}
\usepackage[scr=rsfs,cal=euler]{mathalfa}
\usepackage{xfrac}
\usepackage[unicode,hidelinks,bookmarks=false]{hyperref}
\newcommand{\ie}{i.e.\ }
\newcommand{\cf}{cf.\ }
\newcommand{\resp}{respectively\ }
\newcommand{\Subsection}{\subsection}
\providecommand{\nobreakdash}{\nobreak\hbox{-}}
\setlength{\emergencystretch}{2em}
\multlinegap0pt
\usepackage{amssymb}
\newtheorem{theorem}{Theorem}
\title{Non-K\"ahler Calabi--Yau manifolds and holomorphic geometric structures}
\author{Indranil Biswas, Sorin Dumitrescu}
\date{Source: arXiv:2511.10139v2 (CC BY 4.0)}
\begin{document}
\maketitle

\noindent\emph{Source and licence.} Non-K\"ahler Calabi--Yau manifolds and holomorphic geometric structures. Version arXiv:2511.10139v2 (CC BY 4.0), distributed by arXiv under the Creative Commons Attribution 4.0 International licence, http://creativecommons.org/licenses/by/4.0/, as recorded in the arXiv licence metadata for this exact version and shown on the abstract page https://arxiv.org/abs/2511.10139v2. Full licence notice: \texttt{licenses/arxiv-2511.10139-CC-BY-4.0.txt}.\par\smallskip

\noindent\textit{Editorial scope.} Counted unit: the article's theorem forms (source0.tex lines 854--857). The article's complete original proof of it is delivered with it (source0.tex lines 859--928, one \texttt{proof} environment of 70 lines). No mathematics is written, repaired or re-derived: the statement and the proof are the article's own lines, in the article's own notation, and no retained line is substituted or re-flowed.  No further source-local context is needed: every cross-reference of every retained line resolves inside the two delivered spans.  Nothing was omitted from any retained span, so the package carries no omission record.  The retained lines cite 1 works, whose entries are transcribed verbatim from the article's own bibliography source in the e-print and delivered with short editorial labels recorded in the item's reference mapping.  No retained line contains a figure environment, an image-inclusion command or drawing code, so no image is delivered and the package declares no asset dependency.  Editorial formatting only, in addition: an AMS \texttt{amsart} preamble that restates the article's own declarations and macros used by the retained lines, namely the environments \texttt{theorem} on separate counters, together with the standard packages those lines need.  The retained environments print in this document as Theorem 1 in order of appearance, whereas the article prints the same items with the numbering of its own full document.  The complete native source of the article as distributed in the arXiv e-print for this exact version is bundled unchanged as \texttt{sources/189018/source0.tex}.
\begin{theorem}\label{simply connected with forms}
There exists a simply connected principal elliptic bundle $M$ over a projective $K3$ surface $B$ admitting 
a (non closed) nonsingular holomorphic one-form $\omega$.
\end{theorem} 

\begin{proof}
Let us first construct a projective $K3$-surface $B$ bearing a holomorphic two-form $\Omega$ such that the
group of periods of $\Omega$ is a lattice $\Lambda$ in $\mathbb C$.

One way to obtain such a pair $(B, \,\Omega)$ is to consider a Kummer surface. Indeed, let $\Gamma$ be the 
standard lattice in $\mathbb{C}^2$ generated by $(1,\,0),\, (\sqrt{-1},\,0),\, (0,\,1),\, (0,\,\sqrt{-1})$. 
Consider the quotient of the abelian variety $Y\,=\,\mathbb {C}^2 / \Gamma$ by the holomorphic involution 
$i\,:\, \mathbb {C}^2 / \Gamma\, \longrightarrow\, \mathbb {C}^2 / \Gamma$ defined by $(x, \,y) 
\,\longmapsto\, (-x,\, -y)$.

The above involution $i$ admits 16 fixed points which are the order 2 points in $Y$. Let $B$ be the 
blow-up of $Y/i $ at these 16 points. The smooth complex projective surface $B$ is a K3 surface (in 
particular, it is simply connected). A holomorphic 2-form $\Omega$ on B is given by the 
pull-back of the translation volume form $dz_1\wedge dz_2$ on $Y$, which is also invariant under the 
involution $i$.

The homology group $H_2(B,\, \mathbb{Z})$ is generated by the 16 copies of ${\mathbb C}P^1$ (the 
exceptional divisors for the above mentioned blow-up) and by the images in $H_2(B,\, \mathbb{Z})$ of the 6 
standard generators of $H_2(Y, \,\mathbb{Z}).$ The integral of $\Omega$ on each of the exceptional divisors is zero 
because the form $\Omega$ is pulled back from $Y/i$ and it hence vanishes on the exceptional divisors. On the 
generators of $H_2(Y,\, \mathbb{Z})$ the integral of $\Omega$ is one of the numbers $0,\, 1/4,\, -1/4,\, 
\sqrt{-1}/4,\, -\sqrt{-1}/4.$ Consequently, the group of periods of $\Omega$ form a lattice $\Lambda$ in 
$\mathbb C$.

Now consider the above pair $(B,\, \Omega)$, with $B$ being the projective K3 surface and $\Omega$ the 
(nonsingular) holomorphic two-form whose group of periods is a lattice $\Lambda$ in $\mathbb C$. The 
corresponding elliptic curve $\mathbb{C}/\Lambda$ will be denoted by $T$.

Let us choose an open cover $\{U_i\}_{i \in I}$ of $B$, by open subsets in the analytic topology, such that 
all $U_i$ and all connected components of $U_i \cap U_j$ are contractible. Then on each open subset $U_i$ 
there exists a holomorphic one-form $\omega_i$ such that $\Omega_i\,=\,d \omega_i$. On nontrivial 
intersections $U_i \cap U_j$, we get that $\omega_i-\omega_j\,=\,dF_{ij}$, with $F_{ij}$ a holomorphic 
function defined on $U_i \cap U_j$.

On intersections $U_i \cap U_j \cap U_k$, the sum $F_{ij} +F_{jk}+F_{ki}$ is a locally constant $2$-cocycle 
that represents the cohomology class in $H^2(B,\, { \mathbb C})$ of the closed two-form $\Omega$. 
Therefore, we can choose the local one-forms $\omega_{i}$ and the associated holomorphic functions $F_{ij}$ 
such that --- for every triple $(i,\, j,\, k)$ --- the value of the $2$-cocycle $F_{ij}+F_{jk}+F_{ki}$ 
lies in the lattice $\Lambda \,\subset\, \mathbb C$ formed by the periods of $\Omega$.

Consider the local holomorphic function $\{F_{ij}\}$ as a $1$-cocycle taking values in the translation 
group of $T$ and form the associated holomorphic principal elliptic bundle $M \, \longrightarrow\, B$ with 
typical fiber $T$ associated to this one-cocycle. The local forms on $U_i \times T$ which are 
$\pi_1^*(\omega_i )+\pi_2^*dz $, where $\pi_1$ (respectively, $\pi_2$) is the projection on the first 
(respectively, second) factor and $d z$ is a translation invariant holomorphic one-form on the elliptic 
curve $T$, glue compatibly to produce a globally defined holomorphic one-form $\omega$ on $M$. By 
construction, the differential $d\omega$ projects on $B$ to $\Omega$ and does not vanish. Hence, the 
holomorphic one-form $\omega$ is not closed; more precisely, its differential $d \omega$ is the pull-back 
on $M$ of the form $\Omega$. Moreover, $\omega \wedge d \omega$ does not vanish at any point in $M$ and 
provides a holomorphic section trivializing the canonical line bundle $K_M$.

Observe that the total space $M$ of the above holomorphic principal bundle
with fiber type $T$ is not K\"ahler because the holomorphic one-form $\omega$ on $M$ is not closed.

Moreover, the above manifold $M$ is simply connected, which is proved in \cite[Theorem 12.3 and Remark 
12.4]{Ho}. More precisely, employing the notation of \cite{Ho} observe that, by construction, the 
topological bundle $M$ is characterized by a characteristic class $c\,=\, a_1 \otimes \lambda_1 + a_2 
\otimes \lambda_2 \,\in\, H^2(M,\, \mathbb Z) \otimes \Lambda$, given by $\lambda_1\,=\,1/4, \lambda_2\,=\, 
\sqrt{-1}/4$ and the evaluation of $a_1$ equals $1$ on the pull-back of the class of the elliptic curve 
$E_1$ in $Y$ defined by the quotient of $\mathbb R (1,0)\oplus \mathbb R (1,0)$, equals $-1$ on the 
pull-back of the class of the elliptic curve $E_2$ in $Y$ defined by the quotient of $\mathbb R 
(0,\,1)\oplus \mathbb R (0,\,1)$ and $a_1$ vanishes on the complement in $H^2(M,\, \mathbb Z)$ of the 
direct summand generated by the pull-back of the classes of $E_1$ and $E_2$. The element $a_2$ is defined 
in a similar way using the pull-back of the classes of the elliptic curves generated by $E_3$ and $E_4$ 
which are the images in $Y$ of $\mathbb R (1,\,0)\oplus \mathbb R (0,\,1)$ and $\mathbb R (0,\,1)\oplus 
\mathbb R (1,\,0)$, respectively. Notice that in our situation the obstruction $\Delta$ in 
\cite[Proposition 6.6 and Remark 12.4]{Ho} vanishes, by construction, and since $a_1,\, a_2$ form the basis 
of a direct summand in $H^2(M,\, \mathbb Z)$ the manifold $M$ is simply connected by \cite[Theorem 
12.3]{Ho}.
\end{proof}

\begin{thebibliography}{99}
\bibitem[Ho]{Ho} T. H\"ofer, Remarks on torus principal bundles, {\it J. Math. Kyoto Univ.}, {\bf 33} (1993), 227--259.
\end{thebibliography}
\end{document}
